脂肪酶催化三羟甲基丙烷（TMP）酯化涉及三个相互竞争的 Ping Pong bi bi 反应，有限的浓度观测难以同时约束底物选择性和动态速率参数。针对两步半机理方法固定指纹选择性、未显式描述伪持续性与批内失活的局限，本文提出带指纹信息先验的伪持续性动力学模型（PPK-FIP）。PPK-FIP 将指纹估计转化为可调宽度的先验，在联合参数估计中引入伪持续性与一阶失活，并通过精确时间重参数化降低模型求值成本。在四项研究的九条反应曲线上，留一时间点交叉验证的合并预测误差较重新拟合的初始简化模型降低 \LotoAllRedOrig\,\%，较初始全模型降低 \LotoAllRedFull\,\%（均为 $p<0.001$）。组件消融表明，预测改善主要来自 Bornadel 数据中的伪持续性修正，以及稀疏采样条件下的指纹先验约束。已知伪持续性真值的模拟中，初始两步法对选择性 $S_M$ 的高估为 \SimSMBiasOrigPct\,\%，PPK-FIP 将对数偏差降至 \SimSMBiasNew。高斯过程残差校正进一步将合并交叉验证误差降低至初始简化模型的约一半，但向前外推中 PPK-FIP 与初始简化模型的合并误差相近。研究结果支持在 TMP 酯化建模中联合评估组成路径、参数约束与样本外预测，并将伪持续性参数解释限定于经过验证的反应条件。
